% Drawing for measurement `mos-breakdown`
% Stands for the deck lines:
%   VA sw da 0 / MA da 0 0 0 DUT / .dc VSW 0 1000 1
%   VB sw gb 0 / MB 0 gb 0 0 DUT / .dc VSW -30 30 0.1
\documentclass[border=8pt]{standalone}
\usepackage[siunitx,RPvoltages,european]{circuitikz}
\ctikzset{bipoles/length=0.9cm, label/align=straight, inductors/scale=0.8}
\ctikzset{inductor=american}
\usepackage{amsmath}
\begin{document}
\begin{circuitikz}[font=\footnotesize, line width=0.5pt]
\draw (0,0) node[nmos](A){};
\node[anchor=west] at ($(A.D)!0.5!(A.S) + (0.55,0)$) {DUT};
\draw (A.S) to[short] ++(0,-0.5) node[ground]{};
\draw (A.D) to[short] ++(0,0.3) to[rmeter, t=A, l={$I_{DSS}$}] ++(0,1.4) to[short] ++(2,0) to[vsource, l={$V_{DS}$ $0\to\SI{1000}{V}$}, invert] ++(0,-2.4) node[ground]{};
\draw (A.G) to[short] ++(-0.8,0) node[ground]{};
\node[anchor=west, align=left] at (-4.6,3.3) {drain--source, gate to source};
\draw (8,0) node[nmos](B){};
\node[anchor=west] at ($(B.D)!0.5!(B.S) + (0.55,0)$) {DUT};
\draw (B.S) to[short] ++(0,-0.5) node[ground]{};
\draw (B.D) to[short] ++(0,0.3) to[short] ++(0,0.3) node[ground, rotate=180]{};
\draw (B.G) to[short] ++(-0.5,0) to[rmeter, t=A, l={$I_{GSS}$}] ++(-1.4,0) to[short] ++(-0.5,0) to[vsource, l_={$V_{GS}$ $\pm\SI{30}{V}$}, invert] ++(0,-1.9) node[ground]{};
\node[anchor=west, align=left] at (6.5,3.3) {gate--source, drain to source};
\end{circuitikz}
\end{document}
